% ============================================================
% 01-introduction.tex — V13 Introduction (condensed)
% ============================================================
\section{Introduction}
\label{sec:intro}

Can Tho University (CTU) serves nearly 49,000 undergraduates in 2025--2026~\cite{ctu2025opening}. Its scale and evolving regulations create advisory needs spanning heterogeneous knowledge forms: (1)~\textbf{academic curricula} encoding hierarchies and prerequisites; (2)~\textbf{tuition schedules} varying by cohort and field; (3)~\textbf{exemption and scholarship policies} combining eligibility conditions with rates; and (4)~\textbf{administrative regulations} comprising lengthy narrative documents. Prior systems address individual domains---CAAS for IT admission counseling~\cite{ma2023caas} and REBot for regulation advising~\cite{ma2025rebot}---but a single representation or retrieval method may be unsuitable across all domains.

RAG and GraphRAG augment LLMs with documents and knowledge graphs~\cite{lewis2020rag,gao2023rag_survey,edge2024graphrag}, while agentic architectures support tool use and multi-agent coordination~\cite{yao2023react,wu2023autogen,schick2023toolformer}. These approaches do not prescribe how heterogeneous institutional evidence should be represented, routed, and processed. Known limitations in language-model arithmetic~\cite{qian2023limitations} further motivate deterministic calculation for financial queries.

\textbf{Central hypothesis.}
When a counseling system must handle diverse knowledge forms and semantically adjacent tools, routing queries to bounded specialists with representation-specific evidence paths may improve reliability relative to a single agent with equivalent capabilities, albeit with coordination latency. Whether this trade-off is favorable requires direct comparative evaluation.

These challenges motivate three gaps: (1)~limited evidence comparing multi-agent and single-agent architectures under matched capabilities; (2)~limited mechanistic understanding of which components drive observed differences; and (3)~limited evidence on how bounded specialization maintains reliability as the tool registry grows. We formulate three research questions:
\begin{itemize}
    \setlength{\itemsep}{1.5pt}
    \setlength{\parsep}{0pt}
    \setlength{\parskip}{0pt}
    \item \textbf{RQ1:} How does the proposed multi-agent architecture compare with a functionally matched single-agent in end-to-end reliability, evidence grounding, and cost?
    \item \textbf{RQ2:} What performance changes are observed when specialist policies, tool isolation, representation-specific evidence access, deterministic computation, and route repair are individually removed or replaced?
    \item \textbf{RQ3:} How does bounded specialist routing affect tool-selection reliability as semantically overlapping tools expand the registry from 11 to 51 tools?
\end{itemize}

This work provides three contributions:
\textbf{C1}---a supervisor-routed bounded multi-agent architecture with four specialists, deterministic route repair, and partitioned tool privileges;
\textbf{C2}---an architecture-level comparative evaluation against a functionally matched single-agent baseline and a routed generic-worker variant; and
\textbf{C3}---architectural ablations of five components and a semantic-neighbor tool-registry stress test across 11--51 tools.

The remainder of this paper is organized as follows. Section~2 reviews related work. Section~3 presents the CTU-Chat architecture. Section~4 describes the experimental setup. Section~5 reports results. Section~6 discusses findings and limitations. Section~7 concludes.
